%Abstract
% The abstract is a self-contained summary of the entire thesis (roughly 150--250 words).
% It must be understandable without reading the rest of the thesis.
% Do not use citations, unexplained abbreviations, or heavy mathematical notation.
%
% Follow the hourglass structure illustrated below: start broad (research area),
% narrow to the specific contribution, then broaden again to implications.

\thispagestyle{empty}
\rule{\textwidth}{\rulethickness}
\begin{center}
\large
\textbf{Abstract}\\
\vspace{0mm} %0mm to get symmetry
\end{center}
\rule{\textwidth}{\rulethickness}
\vspace{1mm}

\textit{Expected visiting times} (EVTs) of the states in Markov Chains are a quantitative measure which are used in a variety of applications. The EVT of a state measures the number of times a model is expected to visit that state. An extension of Markov Chains are interval-valued discrete-time Markov Chains (IDTMCs) for which the exact transition probability function and initial distribution are unknown and instead bounded by intervals.

While EVTs have a broad range of applications and have received increasing attention in the literature, they have not yet been studied for IDTMCs. In this thesis, we examine how EVTs can be computed in IDTMCs.

EVTs of a Markov Chain are closely related to other quantitative measures, such as reachability probabilities and expected rewards, which have already been studied for IDTMCs. We adapt established results for these measures and show how they can be used to calculate EVTs in IDTMCs. Moreover, we prove that EVTs in IDTMCs can be calculated in polynomial time using these approaches and show that the corresponding decision problem is P-complete. Finally, we discuss how EVTs can be used to compute stationary distributions in IDTMCs.